\documentclass[12pt,a4paper]{article}
\usepackage[utf8]{inputenc}
\usepackage[T1]{fontenc}
\usepackage[brazil]{babel}
\usepackage{indentfirst}
\usepackage[margin=2.5cm]{geometry}

\title{Resenha Cr\'itica: Artigo 14 ---\\
Architectural Styles and the Design of Network-based\\
Software Architectures (Fielding, 2000)}
\author{Gustavo Pessoa Firmino Duarte}
\date{23 de Abril de 2026}

\begin{document}
\maketitle

\section{Introdu\c{c}\~ao}

A disserta\c{c}\~ao de doutorado de Roy T. Fielding, defendida em 2000 na
Universidade da Calif\'ornia em Irvine, \'e um dos documentos mais influentes da
hist\'oria da arquitetura de software para sistemas distribu\'idos. Nela, Fielding
define e justifica o estilo arquitetural \textit{Representational State Transfer}
(REST), que se tornaria a base arquitetural da \textit{World Wide Web} e, nas
d\'ecadas seguintes, o modelo dominante para o design de APIs na ind\'ustria de
software. O trabalho n\~ao \'e apenas uma proposta t\'ecnica: \'e uma metodologia
para \textit{avaliar} estilos arquiteturais atrav\'es de suas propriedades e
trade-offs, usando um rigoroso framework de an\'alise que aplica e combina
restri\c{c}\~oes arquiteturais para derivar propriedades de sistema.

\section{REST: Restri\c{c}\~oes e Propriedades Derivadas}

Fielding define REST como um conjunto de seis restri\c{c}\~oes arquiteturais,
cada uma adicionando propriedades ao sistema e potencialmente removendo outras.

A restri\c{c}\~ao \textbf{cliente-servidor} separa as preocupa\c{c}\~oes de
interface do usu\'ario das preocupa\c{c}\~oes de armazenamento de dados, melhorando
a portabilidade e a escalabilidade. A restri\c{c}\~ao \textbf{sem estado}
(\textit{stateless}) exige que cada requisi\c{c}\~ao contenha toda a informa\c{c}\~ao
necess\'aria para ser compreendida --- o servidor n\~ao mant\'em sess\~ao entre
requisi\c{c}\~oes. Isso melhora a visibilidade, a confiabilidade e a escalabilidade,
mas aumenta o volume de dados enviados por requisi\c{c}\~ao.

A restri\c{c}\~ao de \textbf{cache} permite que respostas sejam marcadas como
cacheáveis ou n\~ao, melhorando a efici\^encia e a escalabilidade. A restri\c{c}\~ao
de \textbf{interface uniforme} --- a mais central e distintiva do REST --- simplifica
a arquitetura geral e melhora a visibilidade das intera\c{c}\~oes ao custo de
reduzir a efici\^encia para casos de uso espec\'ificos. Ela se desdobra em quatro
subrestri\c{c}\~oes: identifica\c{c}\~ao de recursos via URIs, manipula\c{c}\~ao
de recursos atrav\'es de representa\c{c}\~oes, mensagens autodescritivas e
\textit{HATEOAS} (hiperm\'idia como motor do estado da aplica\c{c}\~ao).

As restri\c{c}\~oes de \textbf{sistema em camadas} e \textbf{c\'odigo sob demanda}
(opcional) completam o conjunto, permitindo, respectivamente, a interposi\c{c}\~ao
de intermediários transparentes e a transfer\^encia de l\'ogica client-side como
JavaScript.

\section{Conex\~ao com a Pr\'atica Profissional}

A API do meu MVP de e-commerce seguia REST de forma superficial: usava HTTP,
verbos GET/POST/PUT/DELETE e URLs baseadas em recursos (\texttt{/products},
\texttt{/orders}). Mas ao ler Fielding, percebi que a restri\c{c}\~ao mais ignorada
foi o \textit{HATEOAS}: a API n\~ao retornava links para as transi\c{c}\~oes de
estado poss\'iveis, exigindo que o cliente conhecesse de antem\~ao toda a estrutura
da API. Fielding argumenta que \'e exatamente o HATEOAS que permite a evolu\c{c}\~ao
do servidor sem quebrar clientes existentes --- uma propriedade que, reconhe\c{c}o,
teria sido valiosa caso o projeto tivesse evolu\'ido al\'em do MVP.

A restri\c{c}\~ao de aus\^encia de estado tamb\'em me fez refletir sobre a
gest\~ao de sess\~oes de usu\'ario: implementei autentica\c{c}\~ao via JWT
armazenado no cliente, o que \'e RESTful no sentido de que o servidor n\~ao guarda
estado de sess\~ao --- mas havia detalhes de implementa\c{c}\~ao que criavam
depend\^encias impl\'icitas de estado no servidor que violavam essa restri\c{c}\~ao.

Na ArcelorMittal, as integra\c{c}\~oes entre sistemas utilizavam desde APIs REST
modernas at\'e Web Services SOAP legados. A clareza do framework de Fielding \'e
instru\'ivel para avaliar os trade-offs de cada abordagem.

\section{Conclus\~ao}

A disserta\c{c}\~ao de Fielding \'e muito mais do que a defini\c{c}\~ao do REST:
\'e um exemplo exemplar de como avaliar e justificar escolhas arquiteturais atrav\'es
de um framework sistem\'atico de restri\c{c}\~oes e propriedades derivadas. O
REST tornou-se onipresente na ind\'ustria, mas frequentemente de forma superficial,
ignorando restri\c{c}\~oes fundamentais como o \textit{HATEOAS}. A leitura original
de Fielding \'e ess\'encial para compreender n\~ao apenas \textit{o que} \'e REST,
mas \textit{por que} cada restri\c{c}\~ao existe e quais propriedades ela garante
--- um entendimento que permite avaliar cr\'iticamente as escolhas arquiteturais
em qualquer projeto de API.

\end{document}
